\section{Existing System Analysis}
The current system for managing training activities within the organization is largely manual and fragmented. The workflow typically begins with various organizational structures submitting their annual training needs using paper-based documents. These documents include specific details such as the training domain, proposed titles, and lists of potential participants with their professional profiles and departmental affiliations.

Once these needs are collected, the information is manually transcribed into separate Excel workbooks. Each year, a new file is created, which leads to significant challenges in maintaining historical records or tracking the long-term progress of individual employees.

\section{Identified Problems}
The manual process introduces several critical issues that hinder the effectiveness of the training programs:
\begin{itemize}
    \item \textbf{Time Inefficiency}: Repetitive manual data entry and cross-referencing between documents consume substantial administrative resources.
    \item \textbf{Data Inconsistency}: Without automated validation, typographical errors and inconsistent naming conventions lead to duplicate records and unreliable data.
    \item \textbf{Lack of Real-time Tracking}: Monitoring the status of a training session or the attendance of participants requires searching through multiple physical or digital files.
    \item \textbf{Analytical Difficulty}: Generating summary statistics or performance reports is a manual, error-prone task that often yields inaccurate results.
\end{itemize}

\section{Proposed Solution: TrainingMS}
To address these limitations, we propose the development of "TrainingMS", a centralized web application. This solution focuses on the digital transformation of the training workflow through:
\begin{itemize}
    \item \textbf{Centralization}: A unified database to store all information regarding trainers, participants, formations, and historical data.
    \item \textbf{Automation}: Streamlined processes for session planning, participant enrollment, and automated validation of inputs.
    \item \textbf{Reliability}: Enforcement of data integrity through a structured relational database schema.
    \item \textbf{Decision Support}: A dedicated statistics module that provides real-time visualization of training metrics, enabling managers to evaluate the impact of training investments.
\end{itemize}
This transition from manual spreadsheets to a robust web platform ensures a scalable, secure, and efficient management environment.
